\documentclass[11pt,a4paper]{article}
\usepackage{curso}
\begin{document}
\portada{Ejercicios --- Lección}{8.6}{Computadores de Langevin no lineales}
\section{Ejercicios}
\begin{enumerate}[label=\textbf{E\arabic*.}]
\item (B) Media de equilibrio de la neurona cuadrática $U=x^2-Ix$.
\item (B) ¿Por qué el término cuártico satura la respuesta?
\item (B) ¿Qué optimiza «entrenar para un tiempo de lectura»?
\item (B) Nombre dos piezas que faltan antes de afirmar eficiencia energética.
\item (I) Para la neurona cuadrática, calcule $\langle x(t)\rangle$ desde $x(0)=0$ y compárelo con la simulación ($I=2$, $t=1$: $0{,}86$).
\item (I) Para $\beta\to\infty$, ¿cuál es la respuesta de equilibrio de la neurona cuártica con $J=(1,0,1)$? Resuelva $2x+4x^3=I$ para $I=6$.
\item (A) Entrene por algoritmo genético una red de 4 neuronas cuárticas para una función XOR con $t_f=1$.
\end{enumerate}
\section{Soluciones}
\begin{solucion}
\textbf{E1.} Mínimo de $x^2-Ix$ y simetría gaussiana: $\langle x\rangle=I/2$.

\textbf{E2.} El crecimiento $x^4$ domina para $|x|$ grande: grandes entradas apenas desplazan el mínimo ($x\sim(I/4)^{1/3}$).

\textbf{E3.} La salida promedio en $t_f$, sin exigir equilibrio.

\textbf{E4.} Hardware real y mediciones de energía de carga y lectura (y del entrenamiento).

\textbf{E5.} $\dot{\langle x\rangle}=-(2\langle x\rangle-I)\Rightarrow\langle x(t)\rangle=\frac I2(1-e^{-2t})$; $I=2$, $t=1$: $0{,}865$.

\textbf{E6.} $4x^3+2x-6=0\Rightarrow x=1$ (exacto); la media térmica a $\beta=10$ es $0{,}99$.

\textbf{E7.} Con dos neuronas ocultas no lineales, el XOR es aprendible; con solo cuadráticas no.
\end{solucion}
\end{document}
